\section{模型评价与推广}

\subsection{优点}

\begin{enumerate}
    \item \textbf{数据完整性保障：}问题1严格使用仅前3天数据构建预测特征，消除信息泄漏；问题2使用全周期数据进行关联分析（与预测任务形成互补）；问题3综合前两问结论设计策略。三问之间数据窗口分工明确、逻辑递进。
    \item \textbf{因果推断框架：}问题3引入Uplift建模中的四象限分类（Persuadable/Sleeping Dog/Lost Cause），区分推送的增量效应与有机购买，避免了传统响应模型"推给所有人"的效率损失。该框架由Ascarza等\cite{ascarza2025}的RCT实验和Wei等\cite{wei2024}的多干预Uplift网络提供文献支撑。
    \item \textbf{策略验证体系：}构建基线$\to$基准方案$\to$优化方案的三组递进对比，通过5,000次蒙特卡洛模拟验证，优化方案实现67,715元营收（2.8$\times$有机基线），留存率11.4\%满足约束，并通过Bootstrap和敏感性分析量化不确定性。
    \item \textbf{统计完备性：}对关键数值进行了置信区间和假设检验补充：Cox PH诊断识别违背假设的协变量、RSF对照验证模型选择合理性、Bootstrap区间量化钻石阈值和AUC的不确定性、K=2--8多指标比较替代硬编码聚类数。这些补充使模型从"有结果"升级为"有可信度"。
    \item \textbf{早期可部署设计：}问题3通过教师--学生架构（全周期聚类标签$\to$Day1--3 XGBoost分类器，CV准确率92.2\%），使策略分群在新服玩家注册第3天即可实时执行，解决了事后画像与实时部署之间的时间鸿沟。
    \item \textbf{实践可操作：}产出的策略方案包含推送时机、状态触发条件、四象限分类、目标人群、定价和预期贡献，可直接转化为运营配置。
\end{enumerate}

\subsection{不足与改进}

\begin{enumerate}
    \item \textbf{需求曲线不可识别：}由于缺乏"推送但拒绝"的曝光-拒绝数据（A/B测试），本文的需求曲线$f_c(p)=\alpha_c e^{-\beta p}$是假设性近似而非真实需求函数。未来若有随机对照试验数据，可通过Uplift模型（如Two-Model approach或Uplift Random Forest）进行无偏估计。
    \item \textbf{四象限比例依赖假设：}优化方案中Persuadable/Sleeping Dog/Lost Cause的比例源自行业基准（SLG首充转化率10--20\%）和生命周期分布，非从数据直接估计。本文通过$\pm$20\%敏感性分析验证了结论对该假设的稳健性，但更精确的估计需要A/B测试数据。
    \item \textbf{聚类特征含未来信息（已通过教师--学生架构缓解）：}问题3的聚类使用了30天终态特征，本文已通过教师--学生架构（全周期聚类$\to$Day1--3分类器，CV准确率92.2\%）初步解决了实时部署问题。学生模型的准确率虽未达100\%，但在时效性与精度的权衡中是可接受的。
    \item \textbf{样本量有限：}训练数据仅2000名玩家（付费样本75人，3.75\%）。Bootstrap结果表明阶段2 GLM系数的置信区间较宽（如diamond\_median的$\exp(\beta)$ 95\% CI: [0.917, 4.031]），反映了付费子样本稀疏导致的估计不确定性。更大的训练集可改善模型稳定性。
    \item \textbf{Cox比例风险假设部分违背：}PH诊断识别出days\_logged\_d3（$p=0.0001$）和n\_event\_types\_d3（$p=0.014$）违背PH假设——这些特征对流失风险的影响可能随时间递减。本文以Cox为主模型并报告了RSF对照（C-index=0.768 vs Cox 0.826），未来可进一步使用时变系数Cox进行改进。
    \item \textbf{K选择受稳定性约束影响：}K=5基于"最小簇$\geq$1\%"的稳定性约束选出（Silhouette=0.488）。若收紧约束至K=4，聚类粒度变粗、营收降至63,043元但留存降至8.5\%；若放宽至K=6（最小簇0.1\%），营收可达71,045元但微小簇统计不可靠。K=5的选择在细分粒度与统计稳定性之间取得平衡，未来在更大数据集上可重新评估最优K。
    \item \textbf{分布迁移风险：}MC模拟假设新服玩家与历史训练集行为同分布，实际中可能因服务器差异、版本更新等存在分布迁移。建议在新服上线后持续收集数据并更新模型。
\end{enumerate}

\subsection{推广价值}

本文的方法体系具有较强的通用性：
\begin{itemize}
    \item \textbf{SLG/卡牌/RPG类游戏：}用户生命周期特征相似（资源积累$+$社交粘性$+$付费加速），聚类$+$Uplift策略优化框架可直接迁移
    \item \textbf{订阅制产品：}生存分析框架适用于任何订阅制/会员制产品的用户流失分析
    \item \textbf{不同规模验证：}蒙特卡洛模拟可灵活适配任意用户基数，只需调整$N_{\text{sim\_players}}$参数
    \item \textbf{无A/B测试条件下的策略设计：}本文展示的"基准方案+假设性优化方案+敏感性分析"范式，为无随机对照数据时的策略优化提供了可复用的方法论
\end{itemize}
